%%=============================================================================
%% 致谢
%%=============================================================================

\chapter*{致谢}
\addcontentsline{toc}{chapter}{致谢}
\markboth{致谢}{致谢}

时光荏苒，岁月如梭。在厦门大学人工智能系的学习生涯即将画上句号，回首这段求学之路，心中充满感激。在此，谨向所有关心、帮助和支持过我的人致以最诚挚的谢意。

首先，衷心感谢我的导师施明辉副教授。从论文的选题、研究方案的制定，到实验的开展与论文的撰写，施老师都给予了我悉心指导。施老师严谨的治学态度、渊博的学识以及对科研工作的热忱，深深感染和激励着我。在施老师的指导下，我不仅学到了专业知识和研究方法，更懂得了如何做一名合格的科研工作者。施老师的教诲将使我受益终生。

感谢厦门大学以及人工智能专业的每一位老师。感谢你们在课堂上的精彩讲授，让我系统地掌握了人工智能领域的基础理论和专业知识；感谢你们在课后的耐心解答，帮助我攻克了一个又一个学术难题。厦门大学"自强不息，止于至善"的校训精神，已深深融入我的血脉，成为我不断前行的动力源泉。

感谢我的家人。感谢父母多年来含辛茹苦的养育之恩，你们无条件的爱与支持是我求学路上最坚强的后盾。感谢你们在我迷茫时给予鼓励，在我挫折时给予安慰。你们的理解与包容，让我能够心无旁骛地投身学业。

特别感谢当前蓬勃发展的中国人工智能行业，感谢我们伟大的祖国。在本论文的整个研究过程中，DeepSeek 和 Qwen 等国产大语言模型为我提供了极大助力。中国 AI 行业的快速发展，不仅为本研究提供了强有力的技术支撑，更让我深切感受到科技创新的力量与魅力。感谢祖国为我们提供了和平安定的学习环境和日益完善的科研条件。身处新时代，见证着祖国的繁荣发展与科技进步，我深感自豪与幸运。作为新时代的青年，我深知肩上的责任与使命。我将把所学知识奉献给祖国的建设事业，为实现中华民族伟大复兴的中国梦贡献自己的一份力量。

路漫漫其修远兮，吾将上下而求索。毕业不是终点，而是新的起点。带着母校的教诲、师长的期望、家人的支持和祖国的召唤，我将踏上新的征程，不忘初心，砥砺前行。

\begin{flushright}
作者：邓凯璇 \\
\today 于厦门大学
\end{flushright}
